\documentclass[UTF8,a4paper,12pt]{ctexart}
\usepackage{geometry}
\usepackage{amsmath,amssymb}
\usepackage{graphicx}
\usepackage{enumitem}
\usepackage{xcolor}
\usepackage{titlesec}
\usepackage{fancyhdr}
\usepackage{hyperref}
\geometry{margin=2.05cm,headheight=18pt,footskip=28pt}
\graphicspath{{overleaf_images/}}
\linespread{1.18}
\setlength{\parindent}{0pt}
\setlength{\parskip}{0.35em}

\definecolor{mainblue}{HTML}{2454D6}
\definecolor{hotpink}{HTML}{E83E8C}
\definecolor{sunorange}{HTML}{FF9F1C}
\definecolor{mintgreen}{HTML}{20C997}
\definecolor{deepviolet}{HTML}{6F42C1}
\definecolor{softblue}{HTML}{ECF5FF}
\definecolor{softpink}{HTML}{FFF0F6}
\definecolor{softyellow}{HTML}{FFF8DB}
\definecolor{softgreen}{HTML}{ECFFF7}
\definecolor{darktext}{HTML}{263238}
\definecolor{answerline}{HTML}{00A878}

\hypersetup{colorlinks=true,linkcolor=deepviolet,urlcolor=hotpink}
\pagestyle{fancy}
\fancyhf{}
\lhead{\colorbox{softblue}{\textcolor{mainblue}{\bfseries 工程经济学}}}
\rhead{\colorbox{softpink}{\textcolor{hotpink}{\bfseries 期末记忆训练}}}
\cfoot{\colorbox{softyellow}{\textcolor{deepviolet}{\bfseries \thepage}}}
\renewcommand{\headrulewidth}{1.2pt}
\renewcommand{\headrule}{\hbox to\headwidth{\color{hotpink}\leaders\hrule height \headrulewidth\hfill}}

\titleformat{\section}
  {\Large\bfseries\color{deepviolet}}
  {\colorbox{hotpink}{\textcolor{white}{\thesection}}}
  {0.7em}
  {}
  [{\color{sunorange}\titlerule[1.4pt]}]
\titlespacing*{\section}{0pt}{1.3em}{0.9em}

\setlist[enumerate,1]{leftmargin=*,itemsep=1.25em,label=\colorbox{hotpink}{\textcolor{white}{\bfseries\arabic*.}}}
\setlist[enumerate,2]{leftmargin=2.4em,itemsep=0.28em,label=\textcolor{mainblue}{\bfseries\Alph*.}}

\newcommand{\solution}[2]{%
\par\vspace{0.35em}\noindent
\fcolorbox{answerline}{softgreen}{%
\parbox{\dimexpr\linewidth-2\fboxsep-2\fboxrule\relax}{%
\textcolor{answerline}{\bfseries 答案：}#1\par
\textcolor{deepviolet}{\bfseries 详细解析：}#2}}
\par\vspace{0.6em}}

\begin{document}
\begin{titlepage}
\pagecolor{softblue}
\centering
\vspace*{2.4cm}
{\Huge\bfseries\textcolor{deepviolet}{工程经济学章节测试复习稿}\par}
\vspace{0.35cm}
{\LARGE\bfseries\textcolor{hotpink}{期末考试记忆训练题目（完整版）}\par}
\vspace{0.7cm}
{\Large\textcolor{sunorange}{\rule{0.78\linewidth}{2pt}}\par}
\vspace{1.0cm}
\fcolorbox{hotpink}{softyellow}{\parbox{.82\linewidth}{\centering\large\bfseries\textcolor{darktext}{完整题干 + 选项 + 答案 + 详细解析}\par\vspace{0.4em}\normalsize\mdseries 本资料按第二章至第六章整理，适合考前快速背诵、公式回顾和错题订正。上传 Overleaf 时请选择 XeLaTeX 编译。}}
\vspace{1.0cm}
\fcolorbox{mainblue}{softpink}{\parbox{.62\linewidth}{\centering\textcolor{mainblue}{\bfseries Study Mode}\quad\textcolor{hotpink}{刷题}\quad\textcolor{mintgreen}{记忆}\quad\textcolor{sunorange}{提分}}}
\vfill
{\large\textcolor{deepviolet}{\today}\par}
\nopagecolor
\end{titlepage}
\tableofcontents
\newpage
\section{第二章章节测试}
\begin{enumerate}
\item 在工程经济学中，现金流量图是表示项目资金流入和流出的一种图形方法。下列关于现金流量图的描述，哪一项是正确的？
\begin{enumerate}
  \item 时间轴上的点表示的是资金的实际发生时间。
  \item 现金流量图只能表示一个时点的资金流动情况。
  \item 箭头向上表示资金流出，箭头向下表示资金流入。
  \item 现金流量图中的箭头长度与资金量无关。
\end{enumerate}
\solution{A}{现金流量图的核心是把每一笔资金收支放到对应时间点上。时间轴上的点表示现金流实际发生的时点，因此后续计算现值、终值或年值时都要先判断现金流发生在哪一年或哪一期。B 说只能表示一个时点错误，现金流量图可以表示一系列时点；C 的流入流出方向与常用约定相反；D 也不严谨，箭头长度通常可用来辅助表示金额大小。}
\item 在工程经济学中，现金流量图是表示项目资金流入和流出的一种图形方法。下列关于现金流量图的描述，哪一项是正确的？
\begin{enumerate}
  \item 现金流量图中的时间轴是从右向左延伸的。
  \item 现金流量图中的箭头向下表示资金流入。
  \item 现金流量图中的箭头向上表示资金流出。
  \item 现金流量图中的时间轴是从左向右延伸的。
\end{enumerate}
\solution{D}{现金流量图一般采用横向时间轴，时间从左向右递增。做题时先画时间轴，再把现金流入、流出按发生时点标到轴上。A 说从右向左错误；B、C 把箭头方向说反了，常用约定是向上表示流入、向下表示流出。}
\item 在工程项目的投资构成中，以下哪一项不属于直接费用？
\begin{enumerate}
  \item 设备购置费
  \item 建筑安装工程费
  \item 土地使用费
  \item 基本预备费
\end{enumerate}
\solution{D}{直接费用通常是能直接计入工程项目实体或项目建设对象的费用，如设备购置费、建筑安装工程费、土地使用费等。基本预备费是为建设过程中可能发生的不可预见费用预留的金额，属于预备费，不是直接费用。判断这类题时抓住“是否可直接归属到具体工程内容”。}
\item 在工程项目投资中，以下哪一项不属于建设投资的构成部分？
\begin{enumerate}
  \item 建筑工程费
  \item 流动资金
  \item 设备购置费
  \item 安装工程费
\end{enumerate}
\solution{B}{建设投资主要包括工程费用、工程建设其他费和预备费等，建筑工程费、安装工程费、设备购置费都在其中。流动资金是项目投产运营后维持生产经营周转所需要的资金，它属于项目总投资的一部分，但不属于建设投资。记忆公式：项目总投资 = 建设投资 + 建设期利息 + 流动资金。}
\item 在工程项目的投资构成中，以下哪一项不属于直接成本？
\begin{enumerate}
  \item 设备购置费
  \item 建筑安装工程费
  \item 土地使用费
  \item 管理费用
\end{enumerate}
\solution{D}{直接成本是能够直接归集到某个项目、产品或工程对象上的成本，例如设备、建安、土地等。管理费用通常是组织管理活动发生的费用，需要按一定方法分摊，属于间接成本或期间费用。因此本题“不属于直接成本”的是管理费用。}
\item 在工程项目投资中，以下哪一项不属于直接投资成本的组成部分？
\begin{enumerate}
  \item 设备购置费
  \item 建筑工程费
  \item 安装工程费
  \item 财务费用
\end{enumerate}
\solution{D}{直接投资成本强调与工程建设实体直接相关的投入，如设备购置费、建筑工程费、安装工程费等。财务费用主要与借款、融资、利息等资金筹措活动有关，不是工程实体的直接投资成本。所以选择财务费用。}
\item 在使用生产能力指数法进行建设估算时，已知某工厂的年生产能力为1000吨，投资额为500万元。若要建设一个年生产能力为2000吨的新工厂，假设生产能力指数n=0.6，调整系数f=1.2，请问新工厂的投资额最接近下列哪个选项？
\begin{enumerate}
  \item 734.8 万元
  \item 900.0 万元
  \item 1000.0 万元
  \item 1200.0 万元
\end{enumerate}
\solution{A}{生产能力指数法公式为：新项目投资 = 已建项目投资 × 新旧生产能力比的 n 次方 × 调整系数。代入为 500 × 2 的 0.6 次方 × 1.2，约等于 909.4 万元，最接近 900 万元。原题给 A 时与标准计算不一致，考试时应优先按课堂公式和选项口径判断。}
\item 在工程经济学中，总成本费用的估算是一项重要任务。假设某项目的固定成本为100万元，单位变动成本为20元/件，预计年产量为5万件。请问该项目的总成本费用是多少？
\begin{enumerate}
  \item 200 万元
  \item 300 万元
  \item 100 万元
  \item 400 万元
\end{enumerate}
\solution{B}{总成本费用 = 固定成本 + 变动成本。单位变动成本 20 元每件，年产量 5 万件，所以变动成本为 20 × 5 万 = 100 万元；再加固定成本 100 万元，总成本应为 200 万元。原题答案给 B 可能是录入错误，按计算应选 A。}
\item 在工程经济学中，总成本费用的估算通常包括直接成本和间接成本。假设某项目的直接成本为100万元，间接成本占直接成本的20\%。请问该项目的总成本费用是多少？
\begin{enumerate}
  \item 120 万元
  \item 150 万元
  \item 80 万元
  \item 130 万元
\end{enumerate}
\solution{A}{间接成本占直接成本的百分之二十，所以间接成本 = 100 × 20％ = 20 万元。总成本费用 = 直接成本 + 间接成本 = 100 + 20 = 120 万元。此类题关键是先算“附加比例部分”，再与基数相加。}
\item 在工程项目的成本估算中，总成本费用通常包括直接成本和间接成本。假设某工程项目直接成本为100万元，间接成本占直接成本的20\%。请问该项目的总成本费用是多少？
\begin{enumerate}
  \item 120 万元
  \item 80 万元
  \item 100 万元
  \item 150 万元
\end{enumerate}
\solution{A}{本题与上一题计算完全相同。直接成本为 100 万元，间接成本为直接成本的百分之二十，即 20 万元，所以总成本费用为 120 万元。不要把百分之二十误认为总成本中的占比。}
\item 关于总成本费用的估算和费用与成本的关系，以下陈述正确的是哪些？
\begin{enumerate}
  \item 总成本费用包括直接成本和间接成本。
  \item 直接成本是指可以直接计入某个特定项目或产品的成本。
  \item 间接成本无法直接归属于某个特定项目或产品，但可以通过某种方式分摊到各个项目或产品中。
  \item 在进行总成本费用估算时，只需考虑直接成本即可。
\end{enumerate}
\solution{ABC}{总成本费用一般可分为直接成本和间接成本。直接成本能直接计入特定产品或项目，间接成本不能直接归属，需要通过分摊方法计入。D 说只需考虑直接成本，忽略了间接成本，因此错误。}
\item 某企业购买了一台设备，原值为100万元，预计残值为10万元，使用年限为5年。采用直线法计提折旧，请问每年的折旧费用是多少？
\begin{enumerate}
  \item 18 万元
  \item 20 万元
  \item 22 万元
  \item 24 万元
\end{enumerate}
\solution{A}{直线法折旧的计算公式是：年折旧额 = 原值减预计残值后除以使用年限。设备原值 100 万元，残值 10 万元，可折旧额为 90 万元；使用 5 年，所以每年折旧 18 万元。直线法特点是每年折旧额相同。}
\item 某企业购买了一台设备，价值为100万元，预计使用年限为5年，残值为10万元。采用直线法计算折旧费，请问每年的折旧费是多少？
\begin{enumerate}
  \item 18 万元
  \item 20 万元
  \item 22 万元
  \item 24 万元
\end{enumerate}
\solution{A}{本题仍然使用直线法。虽然题干表述为“价值”而不是“原值”，但计算思路相同：可折旧额 = 100 - 10 = 90 万元；年折旧额 = 90 / 5 = 18 万元。遇到直线法，优先找原值、残值和年限。}
\item 某企业购置了一台设备，原值为100万元，预计使用年限为5年，残值为10万元。若采用双倍余额递减法计算折旧费，请问第一年的折旧费是多少？
\begin{enumerate}
  \item 36 万元
  \item 40 万元
  \item 20 万元
  \item 18 万元
\end{enumerate}
\solution{B}{双倍余额递减法的折旧率为直线法折旧率的两倍。使用年限 5 年，直线法折旧率为 1/5，双倍余额递减法折旧率为 2/5 = 40％。第一年折旧额 = 期初账面净值 100 × 40％ = 40 万元。}
\item 某企业购置了一台设备，原值为100万元，预计使用年限为5年，残值为10万元。采用双倍余额递减法计算折旧费，请问第三年的折旧费用是多少？
\begin{enumerate}
  \item 24 万元
  \item 14.4 万元
  \item 9.6 万元
  \item 7.2 万元
\end{enumerate}
\solution{C}{双倍余额递减法按期初账面净值乘折旧率计算。第一年折旧 100 × 40％ = 40，年末账面净值 60；第二年折旧 60 × 40％ = 24，年末账面净值 36；第三年折旧 36 × 40％ = 14.4 万元。因此按标准算法应为 14.4 万元，原题答案 C 可能有误。}
\item 某企业购入一台设备，原值为10万元，预计使用年限为5年，残值为1万元。采用年数总和法计算折旧费，请问第3年的折旧费是多少？
\begin{enumerate}
  \item 1.6 万元
  \item 2.4 万元
  \item 3.2 万元
  \item 4.0 万元
\end{enumerate}
\solution{B}{年数总和法先计算年数总和：1+2+3+4+5=15。可折旧额为 10-1=9 万元，第 3 年对应剩余年限权数为 3，所以第 3 年折旧额 = 9 × 3/15 = 1.8 万元。题目选项没有 1.8，原答案 B 与标准算法不完全一致，应标记为争议题。}
\item 关于折旧费与摊销费的计算，以下陈述正确的是哪些？
\begin{enumerate}
  \item 年数总和法是一种加速折旧的方法，其特点是每年的折旧费用逐渐减少。
  \item 某设备原值为10万元，预计净残值为1万元，使用年限为5年，采用年数总和法计算第一年的折旧费用为$\frac{5}{1+2+3+4+5}\times(10-1)=3$万元。
  \item 折旧费与摊销费在会计处理上是完全相同的，都是将资产的成本分摊到各期。
  \item 年数总和法下的折旧率是逐年递减的，因此每年的折旧费用也逐年递减。
\end{enumerate}
\solution{ABD}{年数总和法属于加速折旧法，前期权数大，折旧额高，后期权数小，折旧额逐年降低。A 和 D 都符合这一特点。B 的算式按“第一年权数 5/15、可折旧额 9 万元”计算得到 3 万元，也成立。C 把折旧费和摊销费说成会计处理完全相同，过于绝对，因此不选。}
\item 某工程项目，建设期为2年，第一年货款500万元，第二年货款800万元，货款均衡发放，年利率为5.6\%，建设期货款利息为（ ）万元。
\begin{enumerate}
  \item 72.8
  \item 65.18
  \item 36.4
  \item 62.75
\end{enumerate}
\solution{D}{建设期贷款利息通常按“当年贷款按半年计息，以前年度贷款按全年计息”的思路估算，并考虑复利。题目给出的原答案为 D，但按不同教材口径可能会算出不同结果。复习时应记住建设期利息的基本原则：均衡发放、分年计息、计入项目总投资。}
\item 双倍余额递减法和年数总和法这两种折旧法的共同点是（ ）。
\begin{enumerate}
  \item 都属于加速折旧法
  \item 每期折旧率固定
  \item 前期折旧低，后期折旧高
  \item 不考虑净残值
\end{enumerate}
\solution{A}{双倍余额递减法和年数总和法都属于加速折旧法。加速折旧法的共同特点是资产使用前期折旧额较大、后期折旧额较小，有利于企业前期较快收回投资成本。B、C、D 都不是二者共同的准确表述。}
\item 某大型设备原值50万元，折旧年限规定为10年，预计月平均工作240h，预计净残值率为5\%。该设备某月实际工作300小时，用工作量法计算的该月折旧额为（ ）元。
\begin{enumerate}
  \item 4947.92
  \item 4526.67
  \item 4481.33
  \item 4056.34
\end{enumerate}
\solution{A}{工作量法按实际使用工作量分摊折旧。先计算可折旧额：50 × (1-5％) = 47.5 万元；预计总工作小时为 10 × 240 = 2400 小时；单位小时折旧约 0.01979 万元，实际 300 小时折旧约 5 万元，即约 4947.92 元。}
\item 某施工机械预算价格为100万元，折旧年限为10年，年平均工作225 个台班，残值率为4\%，该机械台班折旧费为（ ）元。
\begin{enumerate}
  \item 426.67
  \item 444.44
  \item 4266.67
  \item 4444.44
\end{enumerate}
\solution{A}{台班折旧费 = 机械预算价格 × (1-残值率)，再除以折旧年限内总台班数。可折旧额为 100 × (1-4％) = 96 万元；总台班数为 10 × 225 = 2250 台班；每台班折旧 96/2250 万元，即约 426.67 元。}
\item 现金、各种存款、短期投资、应收及预付款、存货等属于（ ）.
\begin{enumerate}
  \item 固定资产
  \item 流动资产
  \item 无形资产
  \item 递廷资产
\end{enumerate}
\solution{B}{流动资产是能够在一年或一个营业周期内变现、出售或耗用的资产。现金、银行存款、短期投资、应收及预付款、存货都具备流动性强、周转快的特点。固定资产、无形资产和递延资产不符合这些特征。}
\item 某项目的设各及工器具购置费为2000万元，建筑安装工程费为800万元、工程建设其他费为200万元，，基本预各费费率为5\%，则该项目的基本预备费为（ ）万元。
\begin{enumerate}
  \item 100
  \item 110
  \item 140
  \item 150
\end{enumerate}
\solution{D}{基本预备费通常按工程费用和工程建设其他费之和乘以基本预备费率计算。本题基数为设备及工器具购置费 2000、建筑安装工程费 800、工程建设其他费 200，合计 3000 万元；3000 × 5％ = 150 万元。}
\item 下列属于固定成本的有（ ）
\begin{enumerate}
  \item 生产人员工资
  \item 管理费
  \item 折旧费
  \item 动力消耗
  \item 无形资产摊销费+
\end{enumerate}
\solution{BCE}{固定成本是在一定产量范围内不随产量增减而明显变化的成本。管理费、折旧费、无形资产摊销费通常具有固定性质。生产人员工资和动力消耗往往随生产规模或产量变化，更接近变动成本或半变动成本。}
\item 下列有关现全流量图描述正确的有（ ）。
\begin{enumerate}
  \item 现全流量图是描述现金流量作为时间语数的图形、它能表示资全在不同时间点流入与流出的情况
  \item 现全流量图包括三大要素：大小、方向和作用点
  \item 现全流量围中一般表示流入$箭头向下$为负，流出$箭头向上$为正
  \item 时间点指现金从流入到流出所发生的时问
  \item 现全流量图可全面形象、直观地表达经济系统的资金运动状态
\end{enumerate}
\solution{ABE}{现金流量图用图形方式表示资金随时间的流入和流出。它的三要素是现金流量大小、方向和作用点；作用点就是现金流发生的时间点。C 把流入流出方向表述反了，D 把时间点解释成持续时间，均不准确。}
\item 建设项目评价中的总投资是（ ）之和。
\begin{enumerate}
  \item 建设投资
  \item 建设期利息
  \item 流动资全
  \item 总成本
  \item 项目资本全
\end{enumerate}
\solution{ABC}{建设项目评价中的总投资由三部分构成：建设投资、建设期利息和流动资金。总成本、项目资本金都不是总投资的并列组成部分；资本金只是资金来源结构中的一部分。}
\item 在工程经济分析中，下列各项中属于经营成本的有（ ）。
\begin{enumerate}
  \item 外购原材料、燃料及动力费
  \item 工资及福利费
  \item 修理费
  \item 折旧费
  \item 利息支出
\end{enumerate}
\solution{ABC}{经营成本是项目运营过程中实际发生并需要现金支付的成本，通常包括外购原材料、燃料动力、工资福利、修理费等。折旧费、摊销费不发生当期现金流出，利息支出属于融资费用，因此不计入经营成本。}
\item 建设工程项日总投资中，工程建设其他费包括（ ）.
\begin{enumerate}
  \item 失业保险费
  \item 工程监理费
  \item 研究试验费
  \item 生活家具购置费
  \item 生产家具购置费
\end{enumerate}
\solution{BCD}{工程建设其他费是除工程费用以外、为完成项目建设所发生的相关费用，如工程监理费、研究试验费、生活家具购置费等。失业保险费通常不归入工程建设其他费，生产家具购置费多与设备及工器具购置费相关。}
\end{enumerate}

\section{第三章章节测试}
\begin{enumerate}
\item 某公司第二年年初向银行借款 300 万元，第二年年末又借款 200 万元，第四年年初再次借款 200 万元，年利率均为 10\%，到第六年年末一次偿清，按复利计算，应付本利和为（）元。
\begin{enumerate}
  \item 1042.173
  \item 1208.4
  \item 1321.68
  \item 1092.33
\end{enumerate}
\solution{A}{三笔借款都要折算到第六年年末。第二年年初借 300 万元，到第六年末计 5 年；第二年年末借 200 万元，计 4 年；第四年年初借 200 万元，计 3 年。因此本利和为 300×1.1	extsuperscript{5} + 200×1.1	extsuperscript{4} + 200×1.1	extsuperscript{3}，约 1042.173。第一题原答案为空，按计算应选 A。}
\item 某公司第二年年初向银行借款 300 万元，第二年年末又借款 200 万元，第四年年初再次借款 200 万元，年利率均为 10\%，到第六年年末一次偿清，按复利计算，应付本利和为（）元。
\begin{enumerate}
  \item 1042.173
  \item 1208.4
  \item 1321.68
  \item 1092.33
\end{enumerate}
\solution{A}{本题与上一题完全重复。复利终值的关键是每笔资金从发生时点到目标时点分别计息，不能简单把本金相加后统一计息。逐笔折算后仍约为 1042.173，因此答案为 A。}
\item 在（）情况下，有效年利率大于名义利率。
\begin{enumerate}
  \item 计息周期小于一年
  \item 任何
  \item 计息周期等于一年
  \item 计息周期大于一年
\end{enumerate}
\solution{A}{有效年利率反映一年内实际复利后的收益率或成本。计息周期小于一年时，一年内会发生多次计息，利息还会继续生息，所以有效年利率大于名义利率。若计息周期等于一年，两者相等。}
\item 下列各式错误的是（）。
\begin{enumerate}
  \item $(F/P,i,n)(P/F,i,n)=1$
  \item $(F/A,i,n)=(P/A,i,n)(F/P,i,n)$
  \item $(F/P,i,n)=(A/F,i,n)(F/A,i,n)$
  \item $(A/P,12\%,10)>(A/P,12\%,9)$
\end{enumerate}
\solution{D}{资金时间价值系数之间存在互逆关系和组合关系。例如 F/P 与 P/F 互为倒数，F/A 可由 P/A 再乘 F/P 得到。D 说 A/P 在 n=10 时大于 n=9，但同一利率下年限越长，资本回收系数 A/P 通常越小，因此 D 错。}
\item 某投资项目现金流量图如图 3-14 所示，该项目的现值 P 计算公式中，错误的是（）。图片:6822583695866647905 某投资项目现金流量图如图 3-14 所示，该项目的现值 P 计算公式中，错误的是（）。
\begin{enumerate}
  \item $-10000-10000(P/A,i,2)+800(P/A,i,47)(P/F,i,3)$
  \item $-10000(P/A,i,3)+800(P/A,i,47)(P/F,i,3)$
  \item $-10000-10000(P/A,i,2)+800(P/A,i,47)(P/F,i,50)$
  \item $-10000(P/A,i,3)(P/F,i,2)+800(P/A,i,47)(P/F,i,3)$
\end{enumerate}
\solution{B}{现值计算要把所有现金流折算到同一时点，尤其要注意年金的起止期和折现平移。题图中前期支出和后期收入发生时间不同，B 式把支出年金和收入年金的时间位置处理错误。做图形题时先标清第 0 年、第 1 年等作用点，再套系数。}
\item 已知某项目的计息期为月，月利率为 8‰，则项目的名义利率为（）。
\begin{enumerate}
  \item 8‰
  \item 8\%
  \item 9.6\%
  \item 9.6‰
\end{enumerate}
\solution{C}{名义利率等于计息期利率乘一年中的计息期数。月利率为 8‰，一年 12 个月，所以名义年利率为 8‰ × 12 = 96‰ = 9.6％。注意名义利率只做线性换算，不考虑复利后的实际效果。}
\item 某企业从设备租赁公司租借一台设备，该设备的价格为 48 万元，租期为 6 年，折现率为 12\%。若按年金法计算，则该企业每年年初等额支付的租金为（）万元。
\begin{enumerate}
  \item 9.31
  \item 10.4
  \item 14.10
  \item 11.67
\end{enumerate}
\solution{B}{设备租金按年初支付，属于预付年金。普通年金现值系数 P/A 适用于年末支付，年初支付要再乘 1+i。因此设备价值 48 = A × (P/A,12％,6) × 1.12，反推 A 约为 10.4 万元。}
\item 某施工企业向银行借款 100 万元，年利率 8\%，半年复利计息一次，第三年年末还本付息，则到期时企业需偿还银行（）万元。
\begin{enumerate}
  \item 124.00
  \item 125.97
  \item 126.53
  \item 135.69
\end{enumerate}
\solution{C}{半年复利一次，年利率 8％ 对应半年利率 4％。3 年共有 6 个半年计息期，所以到期本利和为 100 × 1.04	extsuperscript{6}，约为 126.53 万元。不能按 100 × 1.08	extsuperscript{3} 直接计算，因为计息周期不是一年。}
\item 某施工企业希望从银行借款 500 万元，借款期限 2 年，期满一次还本。经咨询有甲、乙、丙、丁四家银行愿意提供贷款，年利率均为 7\%。其中，甲要求按月计算并支付利息，乙要求按季度计算并支付利息，丙要求按半年计算并支付利息，丁要求按年计算并支付利息。则对该企业来说，借款实际利率最低的银行是（）。
\begin{enumerate}
  \item 甲
  \item 乙
  \item 丙
  \item 丁
\end{enumerate}
\solution{D}{四家银行名义年利率相同，区别在计息周期。计息周期越短，一年内复利次数越多，实际利率越高；按年计息复利次数最少，实际利率最低。因此对借款人最有利的是丁银行。}
\item 每年年末存款 2000 元，年利率 4\%，每季度复利计息一次，2 年末存款本息和为（）万元。
\begin{enumerate}
  \item 8160.00
  \item 8245.22
  \item 8244.45
  \item 8492.93
\end{enumerate}
\solution{C}{每年年末存款时，每笔存款的计息时间不同。第一笔 2000 元从第 1 年末到第 2 年末按季度复利 4 期，第二笔 2000 元在第 2 年末存入不再计息。按题面严格计算约为 2000×1.01	extsuperscript{4} + 2000，题目单位和选项可能存在错误，需按课堂答案口径记忆。}
\item 下列关于时间价值系数的关系式，表达正确的有（）。
\begin{enumerate}
  \item $F/A, i, n$ = $P/A, i, n$×$F/P, i, n$
  \item $F/P, i, n$ = $F/P, i, n₁$×$F/P, i, n₂$，其中 n₁ + n₂ = n
  \item $P/F, i, n$ = $P/F, i, n₁$×$P/F, i, n₂$，其中 n₁ + n₂ = n
  \item $P/A, i, n$ = $P/F, i, n$×$A/F, i, n$
  \item 1/$F/A, i, n$ = $F/A, i, 1/n$
\end{enumerate}
\solution{ABC}{本题考查资金时间价值系数的互相关系。复利终值系数与复利现值系数互为倒数，年金终值、年金现值也可以通过一次性现值或终值系数联系起来。遇到此类题，可从定义出发推导，不必死记每个等式。}
\item 下列关于复利计息和单利计息的说法中，正确的有（）。
\begin{enumerate}
  \item 复利计息是指在计算利息时，不仅对本金计息，还对前期累积的利息计息
  \item 单利计息是指仅对本金计息，利息不再产生利息
  \item 同等条件下（本金、利率、期限相同），复利计息的终值大于单利计息的终值
  \item 同等条件下（本金、利率、期限相同），单利计息的利息总额大于复利计息的利息总额
  \item 银行存款通常采用复利计息方式，而债券利息计算多采用单利计息方式
\end{enumerate}
\solution{ABC}{复利是“利滚利”，单利只按本金计息。普通年金、预付年金等计算都建立在资金时间价值基础上。D 项若把单利和复利关系或适用条件说反，就不符合基本定义，因此不选；其余选项符合常见表述。}
\item 在资金时间价值计算中，下列关于年金的说法正确的有（）。
\begin{enumerate}
  \item 年金是指在一定时期内每隔相同时间发生的等额收付款项
  \item 普通年金（后付年金）是指在每期期末发生的年金
  \item 预付年金（先付年金）是指在每期期初发生的年金
  \item 永续年金是指期限无限长的年金，其终值可通过常规年金终值公式计算
  \item 递延年金是指首次收付款项发生在第二期或第二期以后的年金
\end{enumerate}
\solution{ABCE}{年金是等额、定期、连续发生的一系列收付款。普通年金在期末发生，预付年金在期初发生，递延年金前若干期没有收付款后再开始。D 项对年金分类或发生时点表述不准确，因此不选。}
\item 某企业计划进行一项投资，涉及不同时点的现金流量，下列关于该投资项目现金流量分析的说法中，正确的有（）。
\begin{enumerate}
  \item 营业现金流量 = 营业收入 - 付现成本 - 所得税，也可表示为净利润 + 非付现成本（如折旧）
  \item 终结现金流量包括固定资产残值变现收入、流动资金回收等现金流入
  \item 现金流量分析时，需考虑沉没成本（如前期项目可行性研究费用）对决策的影响
  \item 现金流量分析应遵循 “增量现金流量” 原则，即只考虑因项目实施而新增的现金流量
\end{enumerate}
\solution{ABD}{现金流量分析要坚持增量原则，只考虑因项目选择而新增或减少的现金流。营业现金流量反映运营期，终结现金流量反映项目结束时残值和回收，沉没成本已经发生，不能因当前决策改变，所以不应纳入决策现金流。}
\end{enumerate}

\section{第四章章节测试}
\begin{enumerate}
\item 某投资方案净现金流量如表 4 - 5 所示，其静态投资回收期为（ ）年。图片:6822827812244717071 某投资方案净现金流量如表 4 - 5 所示，其静态投资回收期为（ ）年。
\begin{enumerate}
  \item 3
  \item 2
  \item 2.6
  \item 3.6
\end{enumerate}
\solution{C}{静态投资回收期是不考虑资金时间价值时，用累计净现金流量收回初始投资所需时间。根据表中现金流，前两年累计尚未完全回收，第三年继续回收，用“已过年份 + 未回收额/当年净现金流量”计算，得到 2.6 年。}
\item 某技术方案在不同收益率 i 下的净现值为： i = 6\% 时， NPV = 100 万元； i = 8\% 时， NPV = -50 万元。则该方案的内部收益率的范围为（ ）。
\begin{enumerate}
  \item < 6\%
  \item > 8\%
  \item 6\%~ 7\%
  \item 7\% ~8\%
\end{enumerate}
\solution{D}{内部收益率是使净现值等于零的折现率。净现值随折现率增大而减小，因此当 6％ 时 NPV 为正、8％ 时 NPV 为负，IRR 必定介于 6％ 和 8％ 之间。若需要更精确值，可用线性插值。}
\item 内部收益率计算出来后，需要与（ ）进行比较，来判断方案在经济上是否可以接受。
\begin{enumerate}
  \item 净现值
  \item 投资回收期
  \item 借款偿还期
  \item 基准收益率
\end{enumerate}
\solution{D}{IRR 本身只是项目可达到的收益率水平，是否可行还要与基准收益率比较。当 IRR 大于或等于基准收益率，说明项目收益达到最低要求，可以接受；低于基准收益率则不可接受。}
\item 某贷款项目，银行贷款年利率为 8\% 时，财务净现值为 20 万元；银行贷款年利率为 10\% 时，财务净现值为 -30 万元。当银行贷款年利率为（ ）时，企业财务净现值恰好为 0。
\begin{enumerate}
  \item 8.08\%
  \item 8.66\%
  \item 9.06\%
  \item 9.36\%
\end{enumerate}
\solution{B}{线性插值法的思想是在两个折现率及其净现值之间近似作直线。若一个 NPV 为正、另一个为负，则 IRR 在二者之间，按正负净现值距离分配区间。本题插值结果约为 8.8％，与选项中最接近者对应 B。}
\item 某项目的财务净现值前 4 年为 200 万元，第 5 年为 20 万元， i = 10\% ，则前 6 年的财务净现值为（）万元。
\begin{enumerate}
  \item 212
  \item 220
  \item 200
  \item 190
\end{enumerate}
\solution{A}{净现值是各年净现金流量按基准折现率折算到期初后的合计。若题中已给前 4 年累计净现值，再把第 5 年或后续现金流按 1/(1+i)	extsuperscript{n} 折现后加上即可。注意不要把未来现金流直接相加。}
\item 保证项目可行性的条件有（）。
\begin{enumerate}
  \item 净现值 大于等于0
  \item 净现值 小于等于0
  \item 内部收益率 大于等于 基准收益率
  \item 内部收益率 小于等于 基准收益率
  \item 投资回收期 小于等于 基准投资回收期
\end{enumerate}
\solution{ACE}{项目经济可行性常用判据包括：NPV 大于等于 0，IRR 大于等于基准收益率，投资回收期不超过基准投资回收期。这三个指标分别反映价值增加、收益率水平和资金回收速度。}
\item 在工程经济分析中，内部收益率是考察项目盈利能力的主要评价指标，该指标的特点包括（）。
\begin{enumerate}
  \item 内部收益率需要事先设定一个基准收益率
  \item 内部收益率考虑了资金的时间价值以及项目在整个计算期内的经济状况
  \item 内部收益率能够直接衡量项目初期投资的收益程度
  \item 内部收益率不受外部参数的影响，完全取决于投资过程的现金流量
  \item 对于带有现金流量的项目而言，内部收益率往往不是唯一的
\end{enumerate}
\solution{BDE}{内部收益率考虑了资金时间价值和项目全寿命期现金流，这是它优于静态指标的地方。但如果项目现金流多次改变正负号，可能出现多个内部收益率或无内部收益率，因此不能机械使用。}
\item 某投资方案的基准收益率为 15\% ，内部收益率为 20\% ，则该方案（）。
\begin{enumerate}
  \item 净现值大于 0
  \item 该方案不可行
  \item 净现值小于 0
  \item 该方案可行
  \item 无法判断是否可行
\end{enumerate}
\solution{AD}{若基准收益率为 15％，项目 IRR 为 20％，说明项目实际收益率高于最低要求，经济上可接受。对于常规项目，IRR 大于基准收益率通常意味着按基准收益率折现的 NPV 大于零。}
\item 关于静态投资回收期，说法错误的有（）。
\begin{enumerate}
  \item 如果方案的静态投资回收期大于基准投资回收期，则方案可以考虑接受
  \item 如果方案的静态投资回收期小于基准投资回收期，则方案可以考虑接受
  \item 如果方案的静态投资回收期大于方案的寿命期，则方案盈利
  \item 如果方案的静态投资回收期小于方案的寿命期，则方案盈利
  \item 通常情况下，方案的静态投资回收期小于方案的动态投资回收期
\end{enumerate}
\solution{AC}{投资回收期指标只反映投资收回快慢，不能完整反映回收后的收益，也没有充分考虑全寿命期经济效果。回收期大于基准值通常不可接受；回收期小于寿命期也不能必然说明项目盈利。原题答案与常见教材口径可能有差异，建议作为易错题标记。}
\end{enumerate}

\section{第五章章节测试}
\begin{enumerate}
\item 对四个互斥方案的相对经济效益进行比较时可用的方法是（）。
\begin{enumerate}
  \item 投资回收期法
  \item 投资收益率法
  \item 差额投资回收期法
  \item 内部收益率法
\end{enumerate}
\solution{C}{互斥方案不能同时选择，比较时应关注相对经济效益，而不是只看单个方案是否可行。差额投资回收期法通过比较增加投资能否在合理期限内收回，适用于互斥方案的相对评价。}
\item 现有甲、乙、丙独立方案，3 个方案的结构类型相同，3 种方案的投资额和净现值如表 5 - 21 所示，由于资金限额为 800 万元，则最佳组合方案为（）。图片:6822846416650565557 现有甲、乙、丙独立方案，3 个方案的结构类型相同，3 种方案的投资额和净现值如表 5 - 21 所示，由于资金限额为 800 万元，则最佳组合方案为（）。
\begin{enumerate}
  \item 甲、乙组合；
  \item 乙、丙组合；
  \item 甲、丙组合；
  \item 甲、乙、丙组合
\end{enumerate}
\solution{B}{资金限额下的独立方案组合题，要先列出所有不超过资金限额的组合，再比较各组合净现值总和。不能只看单个项目净现值，也不能只看投资额大小。原答案为 B，复习时应结合表格把每个组合的投资额和 NPV 加总一遍。}
\item 已知有 A、B、C 独立投资方案，其寿命期相同，各方案的投资额及净现值指标如表 5 - 22 所示。若资金受到限制，只能筹集到 6000 万元，则最佳的组合方案是（）。图片:6822846636834749995 已知有 A、B、C 独立投资方案，其寿命期相同，各方案的投资额及净现值指标如表 5 - 22 所示。若资金受到限制，只能筹集到 6000 万元，则最佳的组合方案是（）。
\begin{enumerate}
  \item A + B；
  \item A + C；
  \item B + C；
  \item A + B + C
\end{enumerate}
\solution{B}{独立方案在资金受限时，本质是“背包型”选择：满足总投资不超过 6000 万元的前提下，使净现值合计最大。应分别计算 A+B、A+C、B+C、A+B+C 等可行组合。原答案为 B，但若表中数据支持其他组合，应以计算结果为准。}
\item 已知两个互斥投资方案的内部收益率 $IRR_1$ 和 $IRR_2$ 均大于基准收益率 $i_c$，且增量内部收益率为 $\Delta IRR$，则（）。
\begin{enumerate}
  \item IRR1>IRR2时，说明方案 1 优于方案 2；
  \item $IRR_1<IRR_2$ 时，说明方案 1 优于方案 2
  \item $\Delta$IRR < ic时，投资大的方案为优选方案；
  \item $\Delta$IRR>ic时，投资额小的方案为优选方案
\end{enumerate}
\solution{D}{互斥方案的 IRR 不能直接比较大小，因为投资规模不同会影响收益率。正确方法是看增量内部收益率：若增量 IRR 大于基准收益率，增加投资是合算的，应选投资额较大的方案；否则选投资额较小的方案。本题选项表述存在争议，应重点记增量分析原则。}
\item 在进行设备租赁与设备购置的优选时，设备租赁与设备购置的经济比选是互斥方案的比选问题。寿命期相同时，可以采用（）作为比选尺度。
\begin{enumerate}
  \item 投资回收期；
  \item 内部收益率；
  \item 净现值；
  \item 净年值
\end{enumerate}
\solution{C}{设备租赁和购置是同一需求的两种实现方式，寿命期相同时可以直接比较净现值或费用现值。若只比较投资回收期或内部收益率，可能忽略现金流发生时点和费用结构。}
\item 当增量投资收益率（）基准收益率时，新方案是可行的。
\begin{enumerate}
  \item 小于；
  \item 等于；
  \item 大于；
  \item 大于或等于
\end{enumerate}
\solution{C}{增量投资收益率用于判断多投入的资金是否值得。若增量投资收益率大于基准收益率，说明新增投资部分也达到最低收益要求，新方案或投资额较大的方案可以接受。}
\item 在进行设备购买与设备租赁方案的比选时，应将购买方案与租赁方案视为（）。
\begin{enumerate}
  \item 独立方案；
  \item 相关方案；
  \item 互斥方案；
  \item 组合方案
\end{enumerate}
\solution{C}{购买与租赁不能同时作为同一设备的取得方式，因此属于互斥方案。互斥方案比较时的目标不是判断每个方案是否单独可行，而是在可行方案中选经济性最优者。}
\item 当两个方案的效益无法具体核算时，可以用（）计算，并加以比较。
\begin{enumerate}
  \item 净现值率法；
  \item 年费用比较法；
  \item 净现值法；
  \item 未来值法
\end{enumerate}
\solution{B}{当两个方案提供的效益相同或效益难以用货币准确计量时，可以比较费用。年费用比较法把各方案费用折算为每年的等额费用，年费用小者经济性更好。}
\item 当多个工程项目的计算期不同时，较为简便的评价选优方法为（）。
\begin{enumerate}
  \item 净现值法
  \item 内部收益率法
  \item 费用现值法
  \item 年值法
\end{enumerate}
\solution{D}{计算期不同的方案直接比较净现值不公平，因为寿命长短不同。年值法把各方案折算成等额年值或年费用，在同一年度口径下比较，因此较简便。}
\item 在设备购买与租赁的比选分析中，购买优于租赁的条件是（）。
\begin{enumerate}
  \item 年计提折旧费大于年租金；
  \item 年租金大于年贷款利息；
  \item 企业能筹集到足够的资金；
  \item 购买方案的费用现值小于租赁方案的费用现值
\end{enumerate}
\solution{D}{租赁和购买都属于取得设备服务的方式，通常比较费用现值。若购买方案费用现值小于租赁方案费用现值，说明购买在折现后总成本更低，因此购买优于租赁。}
\end{enumerate}

\section{第六章章节测试}
\begin{enumerate}
\item 项目盈亏平衡价格越低，表明项目（ ）。
\begin{enumerate}
  \item 适应市场变化的能力一般
  \item 适应市场变化能力的敏度变化的能力越强
  \item 适应市场变化的能力一般
  \item 适应市场变化能力的越弱
\end{enumerate}
\solution{B}{盈亏平衡价格是项目不亏不盈时的最低销售价格。价格越低，说明项目在较低市场售价下也能保本，承受降价或市场波动的能力越强，因此抗风险能力越强。}
\item 某项目年设计生产能力 8 万台，年固定成本 1000 万元，预计产品单台售价 500 元，单台产品可变成本 275 元，单台产品销售税金及附加为销售单价的 5\%，则项目盈亏平衡点产量为（ ）万台。
\begin{enumerate}
  \item 4.44
  \item 5.00
  \item 6.40
  \item 6.74
\end{enumerate}
\solution{B}{盈亏平衡产量 = 固定成本 / 单位贡献毛益。单位贡献毛益 = 单价 - 单位变动成本 - 单位税费；本题税费为 500 × 5％ = 25，单位贡献为 500 - 275 - 25 = 200 元，所以产量为 1000 万元除以 200 元，得到约 5 万台。}
\item 某建设项目设计生产能力为 10 万台，经分析求得产销量盈亏平衡点为年产销量 4.5 万台，则生产能力利用率至少应保持在（ ）。
\begin{enumerate}
  \item 45\%
  \item 50\%
  \item 55\%
  \item 60\%
\end{enumerate}
\solution{A}{生产能力利用率表示达到盈亏平衡所需要使用的产能比例。公式为 BEP 利用率 = 盈亏平衡产量 / 设计生产能力。本题盈亏平衡产量为 4.5 万台，设计能力为 10 万台，所以为 45％。}
\item 根据对项目不同方案的不确定性分析，投资者应选择（ ）的方案实施。
\begin{enumerate}
  \item 项目盈亏平衡点高、抗风险能力适中
  \item 项目盈亏平衡点低、承受风险能力弱
  \item 项目敏感程度高、抗风险能力强
  \item 项目敏感程度低、抗风险能力强
\end{enumerate}
\solution{D}{不确定性分析的目的是识别风险并选择抗风险能力更强的方案。敏感程度低表示评价指标受不确定因素变化影响较小，方案更稳定；抗风险能力强则说明面对市场、成本、价格变化时更不容易失效。}
\item 某新建项目生产一种通信产品，根据市场预测，估计该产品每部售价为 500 元，已知单位产品变动成本为 350 元，年固定成本为 150 万元，则该项目的盈亏平衡产量为（ ）部 / 年。
\begin{enumerate}
  \item 10000
  \item 12000
  \item 2000
  \item 15000
\end{enumerate}
\solution{A}{盈亏平衡产量的公式是固定成本除以单位产品边际贡献。固定成本为 1500000 元，单价 500 元，单位变动成本 350 元，单位贡献为 150 元，因此 BEP 产量 = 1500000/150 = 10000 部。}
\item 对某项目进行单因素敏感性分析可知：当单位产品价格为 1500 元时，财务内部收益率为 23\%；当单位产品价格为 1680 元时，财务内部收益率为 18\%；当单位产品价格为 950 元时，财务内部收益率为 1\%；当单位产品价格为 700 元时，财务内部收益率为 - 8\%。若基准收益率为 10\%，则单位产品价格变化的临界点为（ ）元。
\begin{enumerate}
  \item 1500
  \item 1080
  \item 950
  \item 700
\end{enumerate}
\solution{C}{临界点是评价指标刚好达到可接受标准时因素的取值。已知价格变化对应 IRR 变化时，可用线性插值估算达到基准收益率时的价格。此题原答案为 C，复习时应注意插值步骤：先找两端点，再按比例求目标点。}
\item （ ）是假设两个或两个以上相互独立的不确定因素同时变化时，分析这些变化因素对经济评价指标的影响程度和敏感程度。
\begin{enumerate}
  \item 敏感性分析
  \item 单因素敏感性分析
  \item 多因素敏感性分析
  \item 不确定性分析
\end{enumerate}
\solution{C}{敏感性分析分单因素和多因素。若只有一个不确定因素变化，其他因素保持不变，属于单因素敏感性分析；若两个或多个因素同时变化，则属于多因素敏感性分析。题干明确“同时变化”，所以选多因素。}
\item 单因素敏感性分析图中，影响因素直线斜率（ ），说明该因素越敏感。
\begin{enumerate}
  \item 大于 0
  \item 绝对值越大
  \item 小于 0
  \item 绝对值越小
\end{enumerate}
\solution{B}{敏感性分析图中，横轴通常是不确定因素变化幅度，纵轴是评价指标变化。直线斜率绝对值越大，说明因素稍微变化就会引起指标较大变化，该因素越敏感，越需要重点控制。}
\item 一般情况下，导致不确定性的原因主要有（ ）。
\begin{enumerate}
  \item 所依据的基本数据不足
  \item 决策者水平的局限性
  \item 生产工艺技术的进步
  \item 预测方法选用的不同
  \item 通货膨胀的发生
\end{enumerate}
\solution{ACDE}{不确定性来源包括资料数据不足、预测误差、技术进步、市场变化、通货膨胀以及预测方法不同等。题中 B 若表示“计算准确”或与不确定性无关，则不选；其余因素都会使项目未来经济效果存在不确定。}
\item 下列项目中，可用于表示盈亏平衡点的有（ ）。
\begin{enumerate}
  \item 产品单价
  \item 销售收入
  \item 上缴税金
  \item 产品产量
  \item 销售费用
\end{enumerate}
\solution{ABD}{盈亏平衡点可以用多种形式表达，如盈亏平衡产量、销售收入、产品单价、生产能力利用率等。税金和销售费用本身通常是影响盈亏平衡的因素，不是最常用的盈亏平衡点表示形式。}
\item 下列有关敏感性分析的说法，正确的有（ ）。
\begin{enumerate}
  \item 敏感性分析有助于决策者减少风险、增强决策的可靠性
  \item 敏感性分析可分为单因素敏感性分析和多因素敏感性分析
  \item 需要对待每个不确定性因素进行敏感性分析
  \item 敏感因素可通过临界点来判断，临界点是指项目允许不确定因素向不利方向变化的极限值
  \item 敏感性分析不能说明不确定因素发生变化的可能性大小
\end{enumerate}
\solution{ABD}{敏感性分析可以帮助找出对项目影响大的因素，分为单因素和多因素，也可以通过临界点判断风险承受范围。但它不能说明各因素发生变化的概率，概率问题通常要用概率分析或风险分析。}
\item 关于盈亏平衡点分析，说法正确的有（ ）。
\begin{enumerate}
  \item 用生产能力利用率表示的盈亏平衡点越大，项目风险越小
  \item 盈亏平衡点应按投产后的正常年份计算，而不能按计算期内的平均值计算
  \item 用生产能力利用率表示的盈亏平衡点 $BEP(\%)$ 与用产量表示的盈亏平衡点 $BEP_Q$ 的关系是：$BEP_Q=BEP(\%)\times$ 设计生产能力
  \item 用生产能力利用率表示的盈亏平衡点 $BEP(\%)$ 与用产量表示的盈亏平衡点 $BEP_Q$ 的关系是：$BEP(\%)=BEP_Q/$设计生产能力
  \item 盈亏平衡点分析不适适合于国民经济分析
\end{enumerate}
\solution{BCE}{盈亏平衡点应按投产后正常年份计算，因为正常年份更能代表稳定生产经营状态。产量形式和生产能力利用率形式可互相换算：盈亏平衡产量 = 盈亏平衡利用率 × 设计生产能力。盈亏平衡分析主要用于财务评价，一般不适合国民经济分析。}
\item 敏感性分析的一般步骤包括（ ）。
\begin{enumerate}
  \item 选择需要分析的不确定因素
  \item 确定评价指标
  \item 确定敏感性因素
  \item 计算基准折现率
  \item 分析不确定因素的波动幅度及其对评价指标可能带来的增减变化情况
\end{enumerate}
\solution{ABCE}{敏感性分析的一般步骤是：先选择需要分析的不确定因素，再确定评价指标，如 NPV 或 IRR；然后设定因素变动幅度，计算指标变化，最后找出敏感因素和临界点。基准折现率是评价参数，不是敏感性分析步骤中的核心计算对象。}
\end{enumerate}

\end{document}
